% GENERATED FILE. Edit canonical lessons, then run npm run generate:books.
% canonical-source-hash: fnv1a64:751fedb8e170b7ea
% canonical-lessons: TE-C184-venna, TE-C184-paccadi, TE-C184-payasam, TE-C184-dosakaya, TE-C184-bendakaya

\chapter{Butter, Chutney, Sweet milk pudding, Cucumber, Okra}
\label{ch:te-butter-chutney-sweet-milk-pudding-cucumber-okra}
\hlchaptermodality{184}

\begin{chapteropening}
\textbf{By the end of this chapter:} \emph{I can say butter, a chutney, a sweet milk pudding, a cucumber and okra.}

Complete the last lesson of chapter 184: I can say butter, a chutney, a sweet milk pudding, a cucumber and okra.
\end{chapteropening}

\section[\texorpdfstring{\te{వెన్న} (venna)}{వెన్న (venna)}]{\te{వెన్న} (venna) — butter}

\label{lesson:TE-C184-venna}



\begin{quote}
Before the new one: say the Telugu for a lid, then the Telugu for a frying pan.
\end{quote}

\subsection*{You'll want to know: \te{వెన్న}}
\textbf{\te{వెన్న}} — \emph{venna} — \textquotedblleft{}butter\textquotedblright{}.

\begin{grammarlens}[title={What you've built}]
One more word for this chapter.
\end{grammarlens}

\subsection*{Your turn}
\begin{itemize}
\raggedright
  \item \emph{Say it:} \emph{venna}
  \item \emph{Say it:} \emph{venna}, once more
  \item \emph{Say it:} say \emph{venna}
  \item \emph{Recall:} say the Telugu for a lid, then the Telugu for a frying pan, then say \emph{venna} again
\end{itemize}

\begin{tcolorbox}[breakable,colback=teal!4,colframe=teal!35!black,title={Before you move on}]
What does \te{వెన్న} mean? (\textquotedblleft{}Butter\textquotedblright{}.) Say it once more.
\end{tcolorbox}
\section[\texorpdfstring{\te{పచ్చడి} (paccaḍi)}{పచ్చడి (paccaḍi)}]{\te{పచ్చడి} (paccaḍi) — a chutney, a pickle}

\label{lesson:TE-C184-paccadi}



\begin{quote}
Before the new one: say the Telugu for a frying pan, then the Telugu for butter.
\end{quote}

\subsection*{You'll want to know: \te{పచ్చడి}}
\textbf{\te{పచ్చడి}} — \emph{paccaḍi} — \textquotedblleft{}a chutney, a pickle\textquotedblright{}.

\begin{grammarlens}[title={What you've built}]
One more word for this chapter.
\end{grammarlens}

\subsection*{Your turn}
\begin{itemize}
\raggedright
  \item \emph{Say it:} \emph{paccaḍi}
  \item \emph{Say it:} \emph{paccaḍi}, once more
  \item \emph{Say it:} say \emph{paccaḍi}
  \item \emph{Recall:} say the Telugu for a frying pan, then the Telugu for butter, then say \emph{paccaḍi} again
\end{itemize}

\begin{tcolorbox}[breakable,colback=teal!4,colframe=teal!35!black,title={Before you move on}]
What does \te{పచ్చడి} mean? (\textquotedblleft{}A chutney, a pickle\textquotedblright{}.) Say it once more.
\end{tcolorbox}
\section[\texorpdfstring{\te{పాయసం} (pāyasaṁ)}{పాయసం (pāyasaṁ)}]{\te{పాయసం} (pāyasaṁ) — a sweet milk pudding}

\label{lesson:TE-C184-payasam}



\begin{quote}
Before the new one: say the Telugu for butter, then the Telugu for a chutney.
\end{quote}

\subsection*{You'll want to know: \te{పాయసం}}
\textbf{\te{పాయసం}} — \emph{pāyasaṁ} — \textquotedblleft{}a sweet milk pudding\textquotedblright{}.

\begin{grammarlens}[title={What you've built}]
One more word for this chapter.
\end{grammarlens}

\subsection*{Your turn}
\begin{itemize}
\raggedright
  \item \emph{Say it:} \emph{pāyasaṁ}
  \item \emph{Say it:} \emph{pāyasaṁ}, once more
  \item \emph{Say it:} say \emph{pāyasaṁ}
  \item \emph{Recall:} say the Telugu for butter, then the Telugu for a chutney, then say \emph{pāyasaṁ} again
\end{itemize}

\begin{tcolorbox}[breakable,colback=teal!4,colframe=teal!35!black,title={Before you move on}]
What does \te{పాయసం} mean? (\textquotedblleft{}A sweet milk pudding\textquotedblright{}.) Say it once more.
\end{tcolorbox}
\section[\texorpdfstring{\te{దోసకాయ} (dōsakāya)}{దోసకాయ (dōsakāya)}]{\te{దోసకాయ} (dōsakāya) — a cucumber}

\label{lesson:TE-C184-dosakaya}



\begin{quote}
Before the new one: say the Telugu for a chutney, then the Telugu for a sweet milk pudding.
\end{quote}

\subsection*{You'll want to know: \te{దోసకాయ}}
\textbf{\te{దోసకాయ}} — \emph{dōsakāya} — \textquotedblleft{}a cucumber\textquotedblright{}.

\begin{grammarlens}[title={What you've built}]
One more word for this chapter.
\end{grammarlens}

\subsection*{Your turn}
\begin{itemize}
\raggedright
  \item \emph{Say it:} \emph{dōsakāya}
  \item \emph{Say it:} \emph{dōsakāya}, once more
  \item \emph{Say it:} say \emph{dōsakāya}
  \item \emph{Recall:} say the Telugu for a chutney, then the Telugu for a sweet milk pudding, then say \emph{dōsakāya} again
\end{itemize}

\begin{tcolorbox}[breakable,colback=teal!4,colframe=teal!35!black,title={Before you move on}]
What does \te{దోసకాయ} mean? (\textquotedblleft{}A cucumber\textquotedblright{}.) Say it once more.
\end{tcolorbox}
\section[\texorpdfstring{\te{బెండకాయ} (beṇḍakāya)}{బెండకాయ (beṇḍakāya)}]{\te{బెండకాయ} (beṇḍakāya) — okra, lady's finger}

\label{lesson:TE-C184-bendakaya}



\begin{quote}
Before the new one: say the Telugu for a sweet milk pudding, then the Telugu for a cucumber.
\end{quote}

\subsection*{You'll want to know: \te{బెండకాయ}}
\textbf{\te{బెండకాయ}} — \emph{beṇḍakāya} — \textquotedblleft{}okra, lady's finger\textquotedblright{}.

\begin{grammarlens}[title={What you've built}]
One more word for this chapter.
\end{grammarlens}

\subsection*{Your turn}
\begin{itemize}
\raggedright
  \item \emph{Say it:} \emph{beṇḍakāya}
  \item \emph{Say it:} \emph{beṇḍakāya}, once more
  \item \emph{Say it:} say \emph{beṇḍakāya}
  \item \emph{Recall:} say the Telugu for a sweet milk pudding, then the Telugu for a cucumber, then say \emph{beṇḍakāya} again
\end{itemize}

\begin{tcolorbox}[breakable,colback=teal!4,colframe=teal!35!black,title={Before you move on}]
What does \te{బెండకాయ} mean? (\textquotedblleft{}Okra, lady's finger\textquotedblright{}.) Say it once more.
\end{tcolorbox}
